% !TeX root = ../main.tex
\section{学位论文工作计划}
\label{sec:work-plan}

\subsection{前期研究基础与已完成工作}
\label{sec:preliminary-work}
\TODO{如实说明已完成的文献整理、理论准备、资料收集或调查准备等工作，以及能够支持后续研究的条件。}

\subsection{研究进度与时间安排}
\label{sec:schedule}
% 各阶段任务直接纳入时间安排；预期成果在表后用一段话说明。
\begin{table}[htbp]
  \centering
  \caption{研究进度与时间安排}
  \label{tab:research-schedule}
  \begingroup
  \renewcommand{\arraystretch}{1.3}
  \begin{tabularx}{\textwidth}{@{}p{0.24\textwidth}X@{}}
    \toprule
    时间安排 & 研究任务 \\
    \midrule
    \TODO{起止年月} & \TODO{明确研究问题，梳理文献，完善研究设计与开题报告。} \\
    \TODO{起止年月} & \TODO{收集、整理研究材料，开展相关调查或分析工作。} \\
    \TODO{起止年月} & \TODO{展开核心论证，完成学位论文初稿。} \\
    \TODO{起止年月} & \TODO{根据反馈补充材料、修改论文，完成定稿及答辩准备。} \\
    \bottomrule
  \end{tabularx}
  \endgroup
\end{table}

各阶段安排见表\ref{tab:research-schedule}。\TODO{在此简述预期形成的学位论文及主要认识成果。}
